% Chapter 2 - System concepts.
% Sources: docs/vision.md, docs/glossary.md, docs/architecture.md,
% claude/decisiones.md (DEC-005, DEC-008, DEC-009, DEC-011, DEC-012,
% DEC-022, DEC-026).
% Target length: 10-12 pages.
%
% Pedagogical chapter. The reader who reaches here has read the
% Introduction but does not yet know anything specific about UnivOrch.
% Concepts are introduced in the order a new reader needs to acquire
% them, from what each user sees to the internal vocabulary, and ends
% just before the implementation details (which belong to Chapter 5).

\chapter{System concepts}
\label{ch:conceptos}

UnivOrch is a service with several moving parts and a vocabulary of
its own. Before reviewing the related work and the development of
the proof of concept, this chapter introduces the ideas the reader
will encounter throughout the rest of the report. The order is
deliberately pedagogical: we start from what each user sees, then
introduce the underlying domain (hypervisors, VMs and their
states), then the orchestrator's central object (the descriptor),
then the tree that organizes descriptors, and only after that the
mechanisms that hold the tree together (cascade inheritance,
declarative loading, lexical closure of reusable resources). The
chapter closes with a glance at how the same design fits the
role-based access control that v1 leaves as future work.

The aim is that, after reading this chapter, a reader knows what
UnivOrch \emph{is and does} without needing to know how it is built.
The implementation details are presented in \Cref{ch:desarrollo}.

%==================================================================
\section{The product through three lenses}
\label{sec:tres-lentes}

UnivOrch is built for three kinds of user, and each of them sees a
different slice of the system. Understanding what each one sees,
and what each one does \emph{not} see, is the easiest way to grasp
what the system is for.

\paragraph{The end user (a student).}
The end user opens the web interface, logs in, and sees a curated
list of the virtual machines assigned to them. The list is small:
typically the machines for one or two courses. The user can start,
stop and (when offered) reset their own machines. They do not see
the catalogue of base templates, the underlying hypervisors, the
hierarchy of folders or the orchestrator's vocabulary; what they see
is a domain-specific view designed for them, in our case the
metaphor of a virtual \emph{desk} that contains \emph{computers}
(DEC-009). The whole point of UnivOrch from this perspective is
that the student operates without depending on the system
administrator.

\paragraph{The manager (a teacher).}
The manager (in the teaching application, a professor) sees a
larger slice. They see the branch of the system that belongs to
their courses: the folders for their subjects, the descriptors for
each student, the templates they have made available to those
students. They can create folders and descriptors inside their
branch, change properties, deploy or undeploy whole subjects in a
single operation, and audit the history. They do not see the
catalogue of users, the global system configuration, the credentials
of hypervisors or the branches that belong to other teachers.

\paragraph{The administrator.}
The administrator sees everything. They define the hypervisors and
the base templates that everyone else uses, they manage the user
catalogue, they grant or revoke access to other roles, and they
respond to incidents (for example, recovering a descriptor that was
left inconsistent after a failed deployment). The administrator is
also the person who installs and operates the daemon itself.

A single person can, of course, hold different roles in different
parts of the system; the role lives on the folder, not on the user
record. We make this precise in \Cref{sec:rbac-vista} at the end of
the chapter. For now, it is enough to think of three different views
of the same product.

%==================================================================
\section{Hypervisors, virtual machines and their states}
\label{sec:hipervisores}

A \emph{hypervisor} is a piece of software that lets multiple
virtual machines share the same physical hardware. VMware ESXi and
Proxmox VE are two examples that this work targets. Each VM
believes it owns its CPU, memory, disks and network adapters; the
hypervisor presents that illusion and arbitrates access to the real
hardware. A hypervisor lives on a single host (or a small cluster of
hosts) and exposes an API that lets external programs create,
destroy, start and stop VMs.

A \emph{virtual machine}, from the hypervisor's perspective, has
three runtime states relevant to this work:

\begin{itemize}
  \item \texttt{running}: the VM is powered on and executing.
  \item \texttt{stopped}: the VM exists in the hypervisor but is
        powered off.
  \item \texttt{paused}: the VM is suspended in place; its memory
        is preserved and a resume operation can return it to
        \texttt{running}.
\end{itemize}

These three states are well known and shared by every hypervisor.
UnivOrch reports them but does not invent them.

\paragraph{The new axis: provisioned versus deployed.}
The most important concept of this chapter is that UnivOrch
introduces a \emph{second state axis} that does not exist at the
hypervisor level. This second axis is the orchestrator's own view of
the relationship between the descriptor and the real VM, and it has
four values (DEC-022, DEC-032):

\begin{description}
  \item[\texttt{provisioned}] -- the orchestrator has the
        descriptor (definition and identifier) but the
        \emph{hypervisor has no VM yet}. The VM has been described
        but not built. This is the natural initial state of every
        descriptor.
  \item[\texttt{deployed}] -- the descriptor has a matching VM in
        the hypervisor. This is the state from which the three
        runtime states above become meaningful: a deployed VM may
        be running, stopped or paused.
  \item[\texttt{broken}] -- an operation (typically a deploy or
        undeploy) failed in a way that left the orchestrator and
        the hypervisor out of sync. Recovery requires an explicit
        \emph{force-undeploy}.
  \item[\texttt{unreachable}] -- the hypervisor cannot be reached;
        the actual state of the VM is therefore unknown. The
        orchestrator marks this state reactively, when a request
        fails because of a connectivity error.
\end{description}

This second axis is the conceptual key of UnivOrch. It lets a
professor describe an entire subject \emph{before} any VM is built,
keep the descriptions versioned and reproducible, and then either
deploy them all at once or let each student deploy their own machine
on demand. Without the \texttt{provisioned} state, the only way to
describe an environment would be to instantiate it, which is what
forces every traditional setup script to be tied to the moment of
deployment.

A descriptor in \texttt{provisioned} consumes no hypervisor
resources. A descriptor in \texttt{deployed} corresponds to a VM
that exists, takes disk space and may consume memory and CPU. Moving
from \texttt{provisioned} to \texttt{deployed} and back is, from the
user's point of view, the central operation of the system.

\Cref{fig:state-machine} on page~\pageref{fig:state-machine} of
\Cref{ch:desarrollo} shows the formal state machine that ties both
axes together.

%==================================================================
\section{The descriptor: definition plus proxy}
\label{sec:descriptor}

The main object UnivOrch manipulates is the \emph{descriptor}. The
name is borrowed from operating systems: in a Unix-like kernel, a
file descriptor is a small integer that points to an open file
without being the file itself; the application reads and writes
through the descriptor, the kernel takes care of the actual file.
The analogy is direct (DEC-005). In UnivOrch:

\begin{itemize}
  \item The descriptor is the orchestrator's handle to a virtual
        machine.
  \item The descriptor is \emph{not} the virtual machine. The VM
        lives in the hypervisor.
  \item Users interact with the descriptor; the orchestrator takes
        care of talking to the hypervisor on their behalf.
\end{itemize}

A descriptor is two things at once. On the one hand, it carries the
\emph{definition} of the VM it represents: the base template it
should be cloned from, the hardware parameters (CPU, memory, disks),
the target hypervisor and any property the system needs in order to
build the VM. On the other hand, it serves as a \emph{proxy} to the
deployed VM: when the descriptor is \texttt{deployed}, it records
the identifier of the VM in the hypervisor, so the orchestrator
knows where to send operations like \texttt{start} or \texttt{stop}.

The distinction between definition and proxy is what makes the
\texttt{provisioned} state meaningful. A provisioned descriptor
carries the full definition but no proxy; a deployed descriptor
carries both. Moving between the two states is what a deploy or an
undeploy does.

A descriptor lives at a single path inside the descriptor tree, and
that path is its identifier. We turn to the tree next.

%==================================================================
\section{The descriptor tree}
\label{sec:arbol}

UnivOrch organizes descriptors in a tree. Each node of the tree is
either a \emph{folder} (an internal node) or a \emph{descriptor}
(a leaf). The structure is identical to a Unix file system, and the
identifiers follow the same convention: a forward slash separates
levels, the root is \texttt{/}, and a node is identified by the
concatenation of names from the root to the node.

For example, the system that the proof of concept ships as a demo
contains:

\begin{verbatim}
/lab/                              folder, administrator's workspace
/lab/networks/                     folder, one subject
/lab/networks/student01            descriptor, a single student's VM
/lab/networks/student02            descriptor, another student's VM
/lab/networks/student03            descriptor, another student's VM
\end{verbatim}

The tree is the conceptual scaffolding of the whole system. It
gives identity to every node (the path), gives structure to access
control (you have permissions on a branch), gives a natural unit
for batch operations (\emph{deploy the whole subject}), and -- as
we see in the next section -- gives a place for shared
configuration.

The tree also gives a natural vocabulary for the teaching
application. A subject becomes a folder. A student becomes either a
descriptor (when each student has a single VM) or a folder
underneath the subject (when each student has several VMs). The same
tree could be used for a CTF competition or a set of laboratory
experiments by changing only the labels in the domain application.

Folders can hold their own properties (description, the resources
they make available to their descendants, the IP pool they assign
addresses from) and zero or more descriptors. A descriptor cannot
contain other descriptors; descriptors are the leaves of the tree.

%==================================================================
\section{Cascade inheritance, with examples}
\label{sec:herencia}

The reason the tree is useful is that \emph{properties cascade}.
Anything defined at a folder is inherited by its descendants. The
mechanism is the closest the system has to a magic ingredient, so
we present it in two passes: first the rules, then two worked
examples.

\paragraph{The three rules.}
The merge rule between parent and child depends on the type of the
field (DEC-026):

\begin{itemize}
  \item \textbf{Scalar fields} (an integer, a string): the child's
        value, if it exists, overrides the parent's.
  \item \textbf{List fields}: the child's list \emph{accumulates}
        on top of the parent's. An item that appears in the
        parent's list is still there for the child unless the parent
        no longer offers it.
  \item \textbf{Dictionary fields}: the parent's and child's
        dictionaries are merged \emph{recursively}, applying the
        same three rules to each key.
\end{itemize}

There are well-known exceptions where the default does not fit. The
IP pool is one: a pool is a coherent block (network, mask, gateway,
range) and merging it field by field would yield nonsense. For
those cases the engine accepts an explicit override that replaces
the whole field rather than merging it.

\paragraph{Example 1: a simple override.}
Suppose the administrator declares, at the folder \texttt{/lab}, a
default of \texttt{cpu: 2} for every machine that lives below. The
folder \texttt{/lab/networks} adds nothing. The descriptor
\texttt{/lab/networks/student01} does not state a CPU value, so it
inherits \texttt{cpu: 2}. The descriptor
\texttt{/lab/networks/student03} explicitly declares \texttt{cpu: 4}.
Both inherit through \texttt{/lab/networks}; \texttt{student01} ends
up with two CPUs, \texttt{student03} with four.

\paragraph{Example 2: an accumulated list.}
The folder \texttt{/lab/networks} declares an \texttt{import:} list
containing \texttt{[linux-vm]}, meaning ``the template
\texttt{linux-vm} is available to my descendants''. The folder
\texttt{/lab/networks/group-A} adds \texttt{import: [windows-vm]}.
A descriptor inside \texttt{group-A} sees both \texttt{linux-vm} and
\texttt{windows-vm}, because the imports of the two folders have
accumulated as the list rule prescribes.

\paragraph{Implications.}
Cascade inheritance is what makes the system practical at scale.
The administrator declares the catalogue of base templates and the
hypervisor configuration \emph{once}, in a high-level folder, and
every descriptor below inherits them. A subject of a hundred
students is then a hundred near-identical descriptors that share
the bulk of their configuration. The teacher edits the shared
configuration in one place; the change reaches every student in the
next deployment.

%==================================================================
\section{Core engine and domain application}
\label{sec:dos-capas}

UnivOrch is built in two layers (DEC-004). The lower layer is the
\emph{generic engine}: it knows about folders, descriptors, jobs,
inheritance and hypervisors, but it knows nothing about subjects,
students or any specific domain. The upper layer is the
\emph{domain application}: it gives meaning to the generic tree in
terms of one concrete use case.

The proof of concept ships a single domain application -- the
teaching one -- but the engine is general. The same engine could
serve a CTF (where the tree models teams instead of subjects), a
set of conference workshops (where each attendee gets a temporary
VM from a shared template), or a research laboratory (where each
experiment has its own folder with its own machines and is
recycled when the experiment ends). The domain application is, in
all of these cases, a thin client that translates the vocabulary of
the domain (\emph{subject}, \emph{student}, \emph{team},
\emph{experiment}) into the generic operations of the engine
(create folder, create descriptor, deploy, undeploy).

The boundary between the engine and the domain application is the
public REST API of the daemon. The teaching application is, from
the engine's point of view, an HTTP client like any other; the
engine has no special path for it.

The separation matters for two reasons. First, the engine can be
reused without modification by a different domain application
written by a third party. Second, the engine is much easier to
maintain when it stays free of teaching vocabulary; the
specialization happens once, at the boundary, and never leaks
inwards.

%==================================================================
\section{Declarative YAML: describing what you want}
\label{sec:declarativo}

UnivOrch follows a declarative model. The user does not say
``create this folder, then create that descriptor, then assign this
template'' one step at a time. The user describes the desired
state in a YAML file, and the system computes the operations needed
to bring the tree into that state. The model is the same one used by
Terraform and Ansible (\Cref{ch:estado-arte}); the operation that
applies it in UnivOrch is called \texttt{load}.

\paragraph{The load operation.}
\texttt{load} takes a YAML file and a destination path. It reads the
YAML, validates it against the current tree, and creates the new
folders and descriptors at the destination. The destination is any
folder in the tree, not only the root; this is what makes a YAML
file \emph{portable}.

For example, the demo file \texttt{demo/setup.yml} describes a
folder \texttt{lab} that contains a folder \texttt{networks} that
contains three student descriptors. The administrator runs:

\begin{verbatim}
univorch load demo/setup.yml /
\end{verbatim}

and the structure appears under the root. A different
administrator, on a different installation, can run:

\begin{verbatim}
univorch load demo/setup.yml /research/site-B/
\end{verbatim}

and the same structure appears under that other folder. The paths
inside the YAML are relative to the destination; the file is a
reusable template.

\paragraph{What a definition looks like.}
The YAML file is a nested dictionary that follows the structure of
the tree. The demonstration file the proof of concept ships looks
like this:

\begin{verbatim}
kind: definition
version: "1"
folders:
  lab:
    description: "Computer Science Lab"
    define hypervisors:
      mock01: { type: mock }
    define templates:
      linux-vm:
        use hypervisor: mock01
        base_vm: linux-base
    folders:
      networks:
        description: "Computer Networks course"
        import: [linux-vm]
        descriptors:
          student01: { use template: linux-vm, cpu: 2 }
          student02: { use template: linux-vm, cpu: 2 }
          student03: { use template: linux-vm, cpu: 4 }
\end{verbatim}

The structure is straightforward. Folders nest under
\texttt{folders}; descriptors live under \texttt{descriptors}.
Resources that are made available to descendants (hypervisors and
templates) are declared with \texttt{define}. Resources that the
folder \emph{borrows} from its parent are declared with
\texttt{import}. Local properties of a descriptor (here, the
\texttt{cpu} of each student) live next to the references.

The vocabulary of \texttt{define}, \texttt{use} and \texttt{import}
deserves its own section.

%==================================================================
\section{Reusable resources: define, use, import, closure}
\label{sec:define-use-import}

The descriptor tree is more than a hierarchy of names; it is also a
hierarchy of \emph{scopes} in which resources are declared, used and
imported. Three keywords govern the model:

\begin{description}
  \item[\texttt{define X:}] declares a resource at this folder. The
        resource exists from now on inside this folder and below.
        The two resources supported in v1 are \texttt{hypervisors}
        and \texttt{templates}. Future extensions will add
        \texttt{datastores} and \texttt{IP pools}, with the same
        grammar.
  \item[\texttt{use X:}] references an existing resource by name. A
        descriptor that says \texttt{use template: linux-vm}
        inherits the fields of the template; a template that says
        \texttt{use hypervisor: mock01} targets that hypervisor.
        The name is resolved in the scope where the resource was
        defined, not where it is used (we return to this below).
  \item[\texttt{import:}] is the visibility control. A folder that
        imports a name from its parent makes that name visible to
        its descendants. A name that is not imported is invisible
        below, even if the parent defined it. This lets an
        administrator publish a large catalogue of templates while
        keeping each course's branch limited to the templates the
        teacher chose to use.
\end{description}

The interaction between \texttt{define}, \texttt{use} and
\texttt{import} mirrors the interaction between
\texttt{define}, \texttt{use} and \texttt{import} keywords in many
programming languages, which is not a coincidence. The next idea
makes the analogy explicit.

\paragraph{Lexical closure of templates.}
When a template defined in folder $A$ is imported into folder $B$,
and a descriptor inside $B$ uses the template, every reference
inside the template (for example \texttt{use hypervisor: mock01})
is resolved \emph{starting from $A$}, not from $B$. The template
carries with it the scope where it was defined; we say it
\emph{closes over} that scope. This is the same behaviour a closure
has in a programming language: the function definition captures the
environment in which it was written.

The motivation is purely practical. A teacher who imports the
template \texttt{linux-vm} from the administrator's catalogue is not
expected to also import the underlying hypervisor; they should just
have to write \texttt{use template: linux-vm} and have it work. The
hypervisor reference inside the template is the administrator's
business, not the teacher's. Lexical closure makes this just work.

\paragraph{Example.}
The demonstration file above declares \texttt{mock01} (the
hypervisor) and \texttt{linux-vm} (the template) at the folder
\texttt{/lab}. The folder \texttt{/lab/networks} imports only
\texttt{linux-vm}; it does not need to import \texttt{mock01}.
When \texttt{student01} says \texttt{use template: linux-vm}, the
engine merges the template's fields into the descriptor; the
template's own \texttt{use hypervisor: mock01} is then resolved
against \texttt{/lab}, where \texttt{mock01} lives, even though
the descriptor itself lives in \texttt{/lab/networks}.

%==================================================================
\section{Why the same design fits a future RBAC}
\label{sec:rbac-vista}

The cascade inheritance described above is not specific to
configuration fields. It applies, with the same rules, to
\emph{role assignments}. UnivOrch defines three roles (DEC-011):

\begin{description}
  \item[\texttt{superuser}] -- the system administrator. Assigned
        at the root; full visibility and control.
  \item[\texttt{manager}] -- a teacher or equivalent. Assigned at
        the folder of a subject; can manage the descendants of that
        folder.
  \item[\texttt{end\_user}] -- a student or equivalent. Assigned at
        the folder of a workstation or set of workstations; can
        operate the descriptors below.
\end{description}

A role assignment is attached to a folder, not to a user record
(DEC-021). The user record itself only carries identity information
(name, password, e-mail). Whether the user is a teacher in one
branch and a student in another is decided by which folders mention
them and in which role.

Two consequences follow directly from cascade inheritance. First,
an assignment made at a high level reaches every descendant
automatically: declaring a user as \texttt{end\_user} at a folder
gives them access to every workstation inside it. Second, a more
specific assignment lower down can override the inherited one if
needed: a user who is an \texttt{end\_user} of an entire subject
can additionally be made a \texttt{manager} of a single
sub-folder of that subject.

The proof of concept does not implement the full RBAC in code; only
the \texttt{superuser} role is realistic in v1, and the service
assumes it runs in a trusted network. The point of this section is
that the model is consistent with the rest of the design and that
\emph{the engine does not need to grow new mechanisms to support
roles}. The same cascade-inheritance code that resolves
configuration fields will resolve role assignments. When the
implementation arrives, it slots into the existing engine without
rearchitecting it.

%==================================================================
\section{Summary}
\label{sec:resumen-conceptos}

The reader who reaches this point has the conceptual map of
UnivOrch:

\begin{itemize}
  \item Three kinds of user (student, teacher, administrator), each
        with a different view of the system.
  \item A second state axis on top of the hypervisor's runtime
        states: \texttt{provisioned}, \texttt{deployed},
        \texttt{broken}, \texttt{unreachable}.
  \item The descriptor as the orchestrator's handle, carrying both
        the definition of a VM and the proxy to the deployed one.
  \item The descriptor tree as the scaffolding that gives identity,
        structure and inheritance to every node.
  \item Cascade inheritance with three type-driven rules and a
        small set of well-defined exceptions.
  \item Two layers, a generic engine and a domain application, with
        the REST API as the boundary between them.
  \item A declarative YAML model and a \texttt{load} operation that
        accepts any destination in the tree.
  \item Three keywords -- \texttt{define}, \texttt{use},
        \texttt{import} -- that turn the tree into a hierarchy of
        scopes, plus a lexical closure that lets templates travel
        with their definition environment.
  \item A role-based access control whose mechanics fit the same
        cascade and are designed for a future implementation.
\end{itemize}

The following chapters take this conceptual map for granted.
\Cref{ch:estado-arte} compares UnivOrch with the closest existing
systems and identifies where the model is original.
\Cref{ch:metodologia} describes how the proof of concept was built.
\Cref{ch:desarrollo} is the technical core of the report: it
documents the architecture and the implementation in detail.
\Cref{ch:conclusiones} closes the report against the objectives.
